\section{Count Partitions with a Given Difference}
\label{sec:dp-count-partitions-diff}

\problemstatement{Given an array $\textit{arr}$ and an integer $d$,
count the number of ways to partition $\textit{arr}$ into two subsets
$S_1$ and $S_2$ such that $S_1 - S_2 = d$ (where $S_1, S_2$ denote the
sums of the two subsets, and every element belongs to exactly one of
them).}

\begin{patternbox}
Another reduction to an already-solved problem: given
$S_1 - S_2 = d$ and $S_1 + S_2 = \textit{totalSum}$ (every element is in
exactly one subset), solving this small system of two equations gives
$S_1 = (\textit{totalSum}+d)/2$ -- turning the partition-counting
question into exactly Section~\ref{sec:dp-count-subsets-sum-k}'s
``count subsets summing to $K$,'' with $K = (\textit{totalSum}+d)/2$.
\end{patternbox}

\subsection*{Understanding the reduction}

Adding the two equations $S_1-S_2=d$ and $S_1+S_2=\textit{totalSum}$
eliminates $S_2$: $2S_1 = \textit{totalSum}+d$, so $S_1 =
(\textit{totalSum}+d)/2$. If $\textit{totalSum}+d$ is odd, no integer
$S_1$ satisfies this, so the answer is immediately $0$. Otherwise, every
way of choosing a subset summing to exactly this $S_1$ corresponds to
exactly one valid partition (the chosen subset is $S_1$, everything else
is $S_2$, and $S_2 = \textit{totalSum}-S_1$ automatically satisfies
$S_1-S_2=d$) -- so the answer is exactly
Section~\ref{sec:dp-count-subsets-sum-k}'s subset-count, evaluated at
this derived target.

\subsection*{All Four Approaches}

Every stage reuses Section~\ref{sec:dp-count-subsets-sum-k}'s
implementation unchanged, called with $k = (\textit{totalSum}+d)/2$.
Only the space-optimized version, plus the reduction wrapper, is shown.

\begin{lstlisting}
#include <bits/stdc++.h>
using namespace std;

int countSubsetsSpaceOptimized(vector<int>& arr, int k) {
    int n = arr.size();
    vector<int> prev(k + 1, 0);

    prev[0] = (arr[0] == 0) ? 2 : 1;
    if (arr[0] != 0 && arr[0] <= k) prev[arr[0]] = 1;

    for (int index = 1; index < n; index++) {
        vector<int> current(k + 1, 0);
        for (int target = 0; target <= k; target++) {
            int notPick = prev[target];
            int pick = (arr[index] <= target) ? prev[target - arr[index]] : 0;
            current[target] = notPick + pick;
        }
        prev = current;
    }
    return prev[k];
}

int countPartitionsWithDiff(vector<int>& arr, int d) {
    int totalSum = accumulate(arr.begin(), arr.end(), 0);
    if ((totalSum + d) % 2 != 0 || totalSum + d < 0) return 0;

    int s1 = (totalSum + d) / 2;
    return countSubsetsSpaceOptimized(arr, s1);
}

int main() {
    vector<int> arr = {1, 1, 2, 3};
    cout << countPartitionsWithDiff(arr, 1) << endl; // 3
    return 0;
}
\end{lstlisting}

\subsection*{Dry run}

\dryrun{Take $\textit{arr} = [1,1,2,3]$, $d=1$.}

$\textit{totalSum}=7$. $S_1 = (7+1)/2 = 4$. Counting subsets summing to
$4$: $\{1,3\}$ (using the first $1$), $\{1,3\}$ (using the second $1$),
and $\{1,1,2\}$ -- three subsets, giving $\textit{answer}=3$. Checking
one of them: subset $\{1,3\}$ has sum $4$; the complementary subset
$\{1,2\}$ has sum $3$; difference $4-3=1=d$, confirmed.

\begin{complexitybox}
\textbf{Time Complexity.} $O(n \cdot \textit{totalSum})$, inherited
from Section~\ref{sec:dp-count-subsets-sum-k}.
\textbf{Space Complexity.} $O(\textit{totalSum})$ for the rolling row.
\end{complexitybox}

\begin{takeawaybox}
Whenever a problem describes a partition in terms of a \textbf{sum and a
difference} of the two parts, solving the resulting two-equation system
for one part's target sum is almost always the fastest path to reusing
an already-solved subset-counting DP, rather than inventing a new
two-dimensional recurrence that tracks both parts directly.
\end{takeawaybox}
